\pagestyle{mystyle}
\chapter{Implementazione del Sistema RAG}
\label{chap:implementazione}

Il capitolo descrive l'implementazione del sistema RAG integrato in i3Guild.
L'obiettivo non è ripetere le scelte architetturali già motivate nel
Capitolo~\ref{chap:requirements_design}, ma mostrare come tali scelte siano
state tradotte in componenti eseguibili: una pipeline di ingestione per
popolare Qdrant, un microservizio FastAPI per il retrieval e la generazione,
un'integrazione applicativa lato Next.js e un workflow agentico per accompagnare
l'utente nella proposta di nuove Gilde. Quest'ultimo viene documentato come
soluzione sperimentale: è stato implementato ed esplorato, ma non è stato
mantenuto come scelta finale a causa del costo computazionale più elevato e del
vantaggio limitato rispetto a un flusso RAG più classico a parità di risorse
hardware.

La realizzazione è distribuita su tre repository applicativi. Il primo è
\texttt{i3guild}, che contiene l'applicazione web Next.js e il database
PostgreSQL~\cite{postgresql2024}. Il secondo è \texttt{i3guild-embedding-pipeline}, responsabile
dell'estrazione dei dati, della normalizzazione testuale, della generazione
degli embedding e dell'upsert su Qdrant. Il terzo è \texttt{i3guild-rag}, un
microservizio Python che espone le API di retrieval e chat. Questa separazione
riduce l'accoppiamento tra applicazione e componente RAG: i3Guild rimane la
sorgente autoritativa dei dati, mentre Qdrant e Ollama vengono trattati come
servizi specializzati~\cite{qdrant2023,ollama2023}.

\begin{figure}[ht]
  \centering
  \begin{tikzpicture}[
    node distance=1.15cm,
    box/.style={draw, rounded corners=2pt, align=center, minimum width=3.1cm,
      minimum height=0.9cm, font=\small},
    arrow/.style={->, thick}
  ]
    \node[box] (postgres) {PostgreSQL\\\texttt{i3guild}};
    \node[box, right=of postgres] (pipeline) {Embedding pipeline\\normalizzazione e chunking};
    \node[box, right=of pipeline] (qdrant) {Qdrant\\indice vettoriale};
    \node[box, below=of qdrant] (rag) {FastAPI RAG\\retrieval e chat};
    \node[box, left=of rag] (ollama) {Ollama\\embedding e generazione};
    \node[box, left=of ollama] (nextjs) {Next.js\\interfaccia i3Guild};

    \draw[arrow] (postgres) -- (pipeline);
    \draw[arrow] (pipeline) -- (qdrant);
    \draw[arrow] (nextjs) -- (ollama);
    \draw[arrow] (nextjs) -- (rag);
    \draw[arrow] (rag) -- (qdrant);
    \draw[arrow] (rag) -- (ollama);
  \end{tikzpicture}
  \caption{Flusso implementativo end-to-end tra sorgente dati, indicizzazione,
  retrieval e generazione.}
  \label{fig:implementation_end_to_end}
\end{figure}

% -------------------------------------------------------------
\section{Struttura dei repository e responsabilità}
\label{sec:implementation_repositories}
% -------------------------------------------------------------

La suddivisione dei moduli segue il flusso dei dati. PostgreSQL contiene le
Gilde, le epoche e le relazioni con gli utenti; la pipeline di embedding legge
queste informazioni e produce documenti Haystack~\cite{haystack2023}; Qdrant memorizza i vettori e
i metadati; il microservizio RAG interroga l'indice e, quando richiesto, invoca
Ollama per generare la risposta. L'applicazione Next.js consuma il servizio
tramite un client HTTP dedicato~\cite{nextjs2024}.

\begin{table}[ht]
  \centering
  \caption{Componenti implementativi principali.}
  \label{tab:implementation_components}
  \small
  \begin{tabular}{p{3.6cm}p{4.7cm}p{5.2cm}}
    \toprule
    \textbf{Repository / modulo} & \textbf{File principali} & \textbf{Responsabilità} \\
    \midrule
    \texttt{i3guild}
      & \texttt{services/rag.service.ts}
      & Client TypeScript verso il microservizio RAG; gestione dello streaming e
        degradazione controllata in caso di errore. \\
    \addlinespace
    \texttt{-embedding-pipeline}
      & \texttt{pipeline/main.py}, \texttt{pipeline/ingest\_sources.py}
      & Estrazione da PostgreSQL, normalizzazione, chunking, embedding con Ollama
        e scrittura su Qdrant. \\
    \addlinespace
    \texttt{-rag}
      & \texttt{app/main.py}, \texttt{app/rag.py}, \texttt{app/pipeline.py},
        \texttt{app/ollama\_service.py}
      & API FastAPI, pipeline Haystack di query, retrieval dinamico, chiamate a
        Ollama e streaming NDJSON. \\
    \addlinespace
    \texttt{-rag/app/agent}
      & \texttt{orchestrator.py}, \texttt{state.py}, \texttt{models.py}
      & Workflow agentico deterministico esplorato per ricerca, ideazione e
        compilazione progressiva della proposta di Gilda; non adottato nella
        soluzione finale. \\
    \bottomrule
  \end{tabular}
\end{table}

La scelta di mantenere l'ingestione fuori dal servizio di query evita che
operazioni batch e richieste interattive competano nello stesso processo. La
pipeline di embedding può essere eseguita manualmente, schedulata o avviata in
container dedicato; il servizio RAG, invece, rimane ottimizzato per rispondere a
richieste HTTP a bassa latenza. La containerizzazione isola dipendenze e servizi
di supporto~\cite{docker2024}.

% -------------------------------------------------------------
\section{Pipeline di ingestione}
\label{sec:implementation_ingestion}
% -------------------------------------------------------------

La pipeline di ingestione è implementata in
\texttt{i3guild-embedding-pipeline/pipeline/main.py}. Il processo viene avviato
come job: verifica le dipendenze, legge le sorgenti dati, produce documenti
Haystack, genera gli embedding tramite Ollama e scrive i risultati in Qdrant~\cite{haystack2023,ollama2023,qdrant2023}. In
caso di errore in una dipendenza essenziale, ad esempio PostgreSQL o Qdrant non
raggiungibile, il processo termina con codice diverso da zero. Questo comportamento
è utile in ambiente Docker o CI, perché rende immediatamente visibile il fallimento
dell'indicizzazione~\cite{docker2024}.

\subsection{Configurazione e controlli preliminari}
\label{subsec:implementation_ingestion_config}

La configurazione avviene tramite variabili d'ambiente. Le più rilevanti sono
\texttt{DB\_URL}, \texttt{QDRANT\_URL}, \texttt{QDRANT\_COLLECTION\_NAME},
\texttt{OLLAMA\_BASE\_URL}, \texttt{OLLAMA\_EMBEDDING\_MODEL} ed
\texttt{EMBEDDING\_DIMENSIONS}. Il modulo legge anche parametri operativi come
\texttt{RECREATE\_INDEX}, \texttt{INGEST\_MAX\_DOCUMENT\_CHARS} e
\texttt{INGEST\_CHUNK\_OVERLAP\_CHARS}.

\begin{table}[ht]
  \centering
  \caption{Parametri principali della pipeline di ingestione.}
  \label{tab:ingestion_configuration}
  \small
  \begin{tabular}{p{4.1cm}p{4.3cm}p{5cm}}
    \toprule
    \textbf{Variabile} & \textbf{Ambito} & \textbf{Uso nell'implementazione} \\
    \midrule
    \texttt{DB\_URL}
      & PostgreSQL
      & Connessione alla sorgente autoritativa delle Gilde. \\
    \texttt{QDRANT\_URL}
      & Qdrant
      & Endpoint del vector database usato in scrittura. \\
    \texttt{QDRANT\_COLLECTION\_NAME}
      & Qdrant
      & Nome della collection da creare o aggiornare. \\
    \texttt{OLLAMA\_BASE\_URL}
      & Ollama
      & Endpoint locale per generare gli embedding. \\
    \texttt{OLLAMA\_EMBEDDING\_MODEL}
      & Ollama
      & Modello usato da \texttt{OllamaDocumentEmbedder}. \\
    \texttt{EMBEDDING\_DIMENSIONS}
      & Indice vettoriale
      & Dimensione dei vettori salvati in Qdrant. \\
    \texttt{RECREATE\_INDEX}
      & Operatività
      & Ricrea la collection quando serve una reindicizzazione completa. \\
    \texttt{INGEST\_MAX\_DOCUMENT\_CHARS}
      & Chunking
      & Limite oltre il quale un documento viene suddiviso. \\
    \bottomrule
  \end{tabular}
\end{table}

Prima dell'ingestione vengono eseguiti tre controlli. Il primo apre una
connessione a PostgreSQL ed esegue una query minimale. Il secondo verifica la
raggiungibilità di Qdrant e crea la collection se non esiste. Il terzo chiama
l'endpoint \texttt{/api/tags} di Ollama per verificare che il server sia
raggiungibile e che il modello di embedding sia configurato. La collection
Qdrant viene creata con distanza coseno e dimensione pari a
\texttt{EMBEDDING\_DIMENSIONS}; nel caso del modello \texttt{nomic-embed-text}
il valore utilizzato è 768~\cite{nussbaum2024nomic}.

\subsection{Estrazione dei dati da PostgreSQL}
\label{subsec:implementation_postgres_extract}

La funzione \texttt{fetch\_all\_guilds()} estrae le informazioni semantiche
delle Gilde tramite una query SQL con join su \texttt{Epoch},
\texttt{GuildUserRelation} ed \texttt{ExternalUser}. Il risultato include i
campi testuali principali della Gilda, i riferimenti all'epoca, la modalità di
partecipazione, l'indicazione di Gilda individuale e la lista dei partecipanti.
Il listato completo dei campi estratti è riportato in
Appendice~\ref{app:B}, Listato~\ref{lst:postgres_extract_fields}.

La query reale include anche una sottoquery per aggregare i partecipanti,
combinando utenti interni ed esterni. Questo passaggio è importante perché i
nomi dei partecipanti sono utili sia come metadati filtrabili sia come contesto
testuale per le domande dell'utente.

\subsection{Normalizzazione e costruzione dei documenti}
\label{subsec:implementation_document_building}

Ogni riga estratta viene trasformata in uno o più oggetti
\texttt{haystack.Document}. Attraverso la funzione \texttt{create\_documents\_from\_row()}
compone un documento testuale con sezioni marcate, ad esempio
\texttt{[SOMMARIO]}, \texttt{[OBIETTIVI]}, \texttt{[RISCHI]},
\texttt{[RISULTATI RAGGIUNTI]} e \texttt{[PARTECIPANTI]}. I campi HTML vengono
normalizzati con \texttt{BeautifulSoup}~\cite{beautifulsoup2024}: il codice rimuove i tag, conserva il
testo e compatta gli spazi. Per il retrieval serve un testo pulito e stabile, non una resa
editoriale. Uno schema semplificato della costruzione del documento è riportato
in Appendice~\ref{app:B}, Listato~\ref{lst:create_document_simplified}.

La strategia descritta nel Capitolo~\ref{chap:requirements_design} prevedeva un
documento consolidato per Gilda. L'implementazione mantiene questa impostazione
come caso normale, ma introduce una soglia di sicurezza:
\texttt{INGEST\_MAX\_DOCUMENT\_CHARS}. Se il testo supera il limite, la funzione
\texttt{split\_text\_with\_overlap()} lo divide in chunk con sovrapposizione
controllata. In questo modo gli outlier non impediscono l'ingestione e, allo
stesso tempo, i chunk successivi mantengono il riferimento alla Gilda tramite
l'intestazione. La granularità dei chunk condiziona qualità del retrieval e
quantità di contesto disponibile per il modello generativo~\cite{chen2023benchmarking}.

Gli identificativi dei documenti sono deterministici:
\texttt{guild-\{guild\_id\}-chunk-\{index\}}. La scelta rende idempotente la
reindicizzazione: rieseguire la pipeline sulla stessa Gilda produce gli stessi
ID e permette a Qdrant di sovrascrivere i punti esistenti invece di duplicarli.
I metadati associati includono \texttt{guild\_id}, \texttt{guild\_name},
\texttt{epoch\_id}, \texttt{is\_individual}, \texttt{participation\_mode},
\texttt{participants}, \texttt{chunk\_index} e \texttt{chunk\_count}.

\subsection{Sorgenti aggiuntive}
\label{subsec:implementation_extra_sources}

Oltre a PostgreSQL, la pipeline supporta due sorgenti opzionali: file JSON e
scraping delle Gilde pubblicate. Il modulo \texttt{ingest\_sources.py} espone
\texttt{load\_guilds\_from\_json()} e \texttt{stream\_guilds\_from\_scraper()}.
Nel primo caso il file viene letto come lista di dizionari; nel secondo caso il
codice usa il web scraper incluso nel repository per iterare sulle Gilde
pubbliche e, se richiesto, arricchirle con dati estratti dagli articoli collegati.

Questa estensione non cambia il contratto con Qdrant: anche i dati esterni
vengono trasformati in documenti Haystack con ID deterministici, contenuto
testuale e metadati. Per le sorgenti non autoritative la policy di duplicazione
predefinita è più conservativa: PostgreSQL usa \texttt{OVERWRITE}, mentre JSON e
scraper usano \texttt{SKIP} salvo configurazione esplicita. La ragione è
pratica: il database applicativo è la fonte principale, mentre le sorgenti
esterne possono essere incomplete o già presenti nell'indice.

\subsection{Embedding e scrittura in Qdrant}
\label{subsec:implementation_embedding_write}

La fase finale costruisce una pipeline Haystack minimale composta da
\texttt{OllamaDocumentEmbedder} e \texttt{DocumentWriter}. L'embedder usa il
modello configurato in \texttt{OLLAMA\_EMBEDDING\_MODEL}, con
\texttt{num\_ctx=8192}; il writer persiste i documenti nel
\texttt{QdrantDocumentStore}. La configurazione essenziale della pipeline è
riportata in Appendice~\ref{app:B}, Listato~\ref{lst:embedding_pipeline_impl}.

La pipeline non usa un componente di cleaning Haystack separato, perché la
normalizzazione dei campi avviene prima della costruzione dei documenti. Questa
scelta rende più esplicito il formato testuale indicizzato e consente di
controllare quali campi entrano nel corpus.

% -------------------------------------------------------------
\section{Microservizio RAG}
\label{sec:implementation_rag_service}
% -------------------------------------------------------------

Il microservizio RAG è implementato con FastAPI nel repository
\texttt{i3guild-rag}. Espone endpoint per retrieval, chat standard, chat in
streaming e workflow agentico. La configurazione è centralizzata in
\texttt{app/settings.py} tramite \texttt{pydantic-settings}; i valori di default
permettono lo sviluppo locale, mentre gli ambienti Docker possono sovrascriverli
con variabili d'ambiente~\cite{fastapi2024,pydantic2024}.

\begin{table}[ht]
  \centering
  \caption{Endpoint principali esposti dal microservizio RAG.}
  \label{tab:rag_service_endpoints}
  \small
  \begin{tabular}{p{3.6cm}p{2.2cm}p{7.2cm}}
    \toprule
    \textbf{Endpoint} & \textbf{Formato} & \textbf{Responsabilità} \\
    \midrule
    \texttt{POST /rag/retrieve}
      & JSON
      & Recupera documenti semanticamente pertinenti e restituisce metadati di
        diagnostica sul taglio dei risultati. \\
    \texttt{POST /chat/reply}
      & JSON
      & Esegue retrieval, costruisce il prompt aumentato e restituisce una
        risposta completa. \\
    \texttt{POST /chat/stream}
      & NDJSON
      & Espone la chat incrementale, separando eventi RAG, token, warning ed
        errori. \\
    \texttt{POST /chat/agentic}
      & JSON
      & Applica l'orchestrazione deterministica e aggiorna lo stato della
        proposta. \\
    \texttt{POST /chat/agentic/stream}
      & NDJSON
      & Restituisce decisione agentica, risultati RAG e risposta in forma
        incrementale. \\
    \bottomrule
  \end{tabular}
\end{table}

\subsection{Avvio, middleware e gestione errori}
\label{subsec:implementation_fastapi}

Il file \texttt{app/main.py} definisce l'applicazione FastAPI e registra CORS,
rate limiting e handler globale delle eccezioni. Il rate limiter usa
\texttt{slowapi} con limite predefinito di 100 richieste al minuto per IP sugli
endpoint standard e 60 richieste al minuto sugli endpoint agentici. Il sistema
non tratta tutti gli errori allo stesso modo: se il messaggio suggerisce una
dipendenza non raggiungibile, ad esempio Qdrant o Ollama, la risposta HTTP è
503; negli altri casi viene restituito 500.

Per ridurre il rischio operativo in ambienti con risorse limitate, il servizio
controlla la memoria disponibile prima delle chiamate a modelli locali. La
funzione \texttt{get\_model\_memory\_requirement()} stima il requisito RAM in
base al nome del modello; \texttt{get\_ram\_warning()} produce un warning se la
memoria libera è inferiore alla soglia. Il warning non blocca la richiesta, ma
viene propagato nella risposta quando disponibile.

\subsection{Endpoint di retrieval}
\label{subsec:implementation_retrieve_endpoint}

L'endpoint \texttt{POST /rag/retrieve} riceve una query, parametri di retrieval
e filtri opzionali. Il modello Pydantic \texttt{RetrieveRequest} definisce i
valori di default: \texttt{top\_k=20}, \texttt{score\_threshold=0.7},
\texttt{max\_documents=10} e \texttt{gap\_threshold=0.08}. La risposta include
i documenti selezionati, il numero di candidati iniziali, la soglia usata, il
motivo del taglio, l'eventuale query riscritta e gli eventuali filtri estratti
automaticamente. Lo schema della richiesta è riportato in
Appendice~\ref{app:B}, Listato~\ref{lst:retrieve_request}.

L'endpoint è volutamente sottile: valida l'input, esegue il controllo RAM e
delega la logica a \texttt{retrieve\_documents()}. Questa separazione rende più
semplice testare il retrieval senza dover passare da FastAPI.

\subsection{Pipeline Haystack di query}
\label{subsec:implementation_query_pipeline}

La pipeline di query è costruita in \texttt{app/pipeline.py}. Il
\texttt{QdrantDocumentStore} punta alla collection configurata e usa embedding
di dimensione 768 con similarità coseno. Il componente
\texttt{OllamaTextEmbedder} trasforma la query utente in vettore; il
\texttt{QdrantEmbeddingRetriever} esegue la ricerca in
Qdrant~\cite{haystack2023,qdrant2023}. Il listato della pipeline è riportato in
Appendice~\ref{app:B}, Listato~\ref{lst:query_pipeline_impl}.

La pipeline è inizializzata in modo lazy e protetta da un lock. Alla prima
richiesta viene creato l'oggetto Haystack e vengono predisposti gli indici sui
payload Qdrant per i campi usati più spesso nei filtri:
\texttt{epoch\_id}, \texttt{is\_individual}, \texttt{participation\_mode},
\texttt{participants} e \texttt{guild\_name}. Le chiamate successive riusano la
stessa pipeline in memoria. La ricerca vettoriale su Qdrant usa indici
approssimati per nearest neighbor; HNSW è il riferimento adottato per questa
classe di problemi~\cite{malkov2018efficient}.

% -------------------------------------------------------------
\section{Logica di retrieval}
\label{sec:implementation_retrieval_logic}
% -------------------------------------------------------------

La funzione \texttt{retrieve\_documents()} in \texttt{app/rag.py} contiene la
parte più specifica del sistema. Il retrieval non è una semplice chiamata a
Qdrant: prima della ricerca la query viene normalizzata e, quando possibile,
riscritta; dopo la ricerca i risultati vengono deduplicati e filtrati con una
logica dinamica. Questa struttura segue il principio RAG di separare conoscenza
parametrica e recupero esplicito di documenti esterni~\cite{lewis2020retrieval}.

\subsection{Riscrittura della query e filtri}
\label{subsec:implementation_query_rewriting}

La funzione \texttt{extract\_topic\_query()} rimuove parole frequenti nel dominio
i3Guild, come "gilda", "fondare" o "consigli", quando non contribuiscono al tema
ricercato. Se dalla frase "vorrei fondare una gilda sulla musica" viene estratto
il topic "musica", la ricerca viene eseguita su tale topic. La risposta conserva
la query riscritta nel campo \texttt{topic\_query}, così il frontend può
mostrarla o usarla per diagnostica.

In parallelo, \texttt{extract\_filters\_from\_query()} cerca vincoli espressi in
linguaggio naturale: epoca, Gilde individuali o di gruppo, modalità di
partecipazione e partecipanti. Se l'utente passa anche filtri espliciti, il
servizio li combina con quelli estratti tramite un operatore \texttt{AND}. Questo
permette di usare la ricerca vettoriale senza perdere vincoli strutturati
presenti nei metadati.

\subsection{Cache e deduplicazione}
\label{subsec:implementation_cache_dedupe}

Per evitare interrogazioni ripetute identiche, \texttt{retrieve\_documents()}
mantiene una cache in memoria con TTL configurabile tramite
\texttt{rag\_cache\_ttl\_seconds}. La chiave include query effettiva, parametri
di retrieval e filtri. Si tratta di una cache locale al processo, sufficiente per
lo scenario prototipale; in un deployment multi replica andrebbe sostituita con
una cache condivisa o disabilitata per evitare differenze tra istanze.

Dopo il recupero da Qdrant, il servizio rimuove duplicati su tre livelli. Prima
controlla l'ID del documento, poi usa una fingerprint testuale normalizzata e
infine confronta i testi con \texttt{difflib.SequenceMatcher}. La soglia fuzzy è
configurabile con \texttt{rag\_dedup\_fuzzy\_threshold}. Questa parte è stata
aggiunta perché il corpus può contenere documenti equivalenti provenienti da
sorgenti diverse o chunk molto simili generati dalla sovrapposizione.

\subsection{Selezione dinamica dei documenti}
\label{subsec:implementation_dynamic_filtering}

La selezione finale avviene in \texttt{filter\_documents\_dynamically()}. I
documenti vengono ordinati per score, filtrati con una soglia minima e poi
tagliati quando il calo relativo tra due score consecutivi supera
\texttt{gap\_threshold}. Se il numero di risultati resta troppo alto, viene
applicato \texttt{max\_documents}. La risposta indica il motivo del taglio:
\texttt{threshold}, \texttt{relative\_gap}, \texttt{max\_cap} o
\texttt{exhausted}. Uno schema della logica di taglio è riportato in
Appendice~\ref{app:B}, Listato~\ref{lst:dynamic_filtering}.

Questa scelta evita un limite \texttt{top\_k} fisso. Per query molto specifiche
il sistema restituisce pochi documenti; per query più ampie può mantenere più
contesto, entro il cap configurato. Il risultato è più stabile per l'utente e
riduce la probabilità di passare al modello generativo documenti debolmente
pertinenti. Ridurre contesto non pertinente è utile anche per mitigare risposte
non fondate sulle fonti recuperate~\cite{shuster2021retrieval}.

% -------------------------------------------------------------
\section{Generazione con Ollama}
\label{sec:implementation_generation}
% -------------------------------------------------------------

Il modulo \texttt{app/ollama\_service.py} gestisce le chiamate al modello di
chat. Le funzioni principali sono \texttt{generate\_chat\_reply()} per risposte
complete e \texttt{generate\_chat\_reply\_stream()} per risposte in streaming.
Entrambe usano l'endpoint \texttt{/api/chat} di Ollama~\cite{ollama2023}.

\subsection{Risoluzione del modello}
\label{subsec:implementation_model_resolution}

Il modello configurato in \texttt{OLLAMA\_LLM\_MODEL} viene confrontato con la
lista dei modelli installati ottenuta da \texttt{/api/tags}. Il codice accetta
sia corrispondenze esatte sia varianti con tag, ad esempio \texttt{:latest}. Se
la configurazione è ambigua o il modello non è installato, il servizio restituisce
un errore esplicito. Questa verifica evita fallimenti meno leggibili durante la
generazione.

La lista dei modelli disponibili viene mantenuta in cache per pochi secondi
(\texttt{ollama\_model\_cache\_ttl\_seconds}), così il servizio non interroga
Ollama a ogni richiesta. La cache è locale al processo e ha solo funzione di
riduzione dell'overhead.

\subsection{Costruzione del prompt aumentato}
\label{subsec:implementation_augmented_prompt}

Prima di inviare la richiesta a Ollama, \texttt{\_prepare\_messages()} costruisce
la lista dei messaggi. Il system prompt, se presente, viene aggiunto per primo.
La cronologia conversazionale viene limitata agli ultimi
\texttt{max\_history\_messages}, in modo da controllare la dimensione del
contesto. Se il retrieval ha restituito documenti, il testo viene aggiunto al
messaggio utente dentro una sezione esplicita:
\texttt{--- CONTESTO RAG ---}. L'aumento del prompt con documenti recuperati è
il meccanismo con cui RAG rende disponibili al modello informazioni esterne
senza modificare i pesi~\cite{lewis2020retrieval}. Il formato del messaggio
aumentato è riportato in Appendice~\ref{app:B},
Listato~\ref{lst:augmented_prompt}.

Ogni documento viene troncato a \texttt{rag\_max\_chars\_per\_doc} caratteri.
Il troncamento è intenzionale: una singola Gilda molto lunga non deve consumare
tutta la finestra di contesto né impedire al modello di vedere altre fonti.

\subsection{Streaming NDJSON}
\label{subsec:implementation_streaming}

Per la chat interattiva il servizio espone \texttt{POST /chat/stream}. Prima
invia un evento \texttt{rag} con i metadati del retrieval, poi inoltra gli eventi
generati da Ollama. La risposta HTTP usa il media type
\texttt{application/x-ndjson}; ogni riga è un oggetto JSON indipendente. Questo
formato è semplice da consumare lato frontend e permette di distinguere token,
warning, errori e fine stream. Un esempio di sequenza NDJSON è riportato in
Appendice~\ref{app:B}, Listato~\ref{lst:ndjson_standard}.

% -------------------------------------------------------------
\section{Workflow agentico}
\label{sec:implementation_agentic}
% -------------------------------------------------------------

Il sistema ha esplorato anche una modalità agentica, esposta dagli endpoint
\texttt{/chat/agentic} e \texttt{/chat/agentic/stream}. La documentazione
progettuale iniziale prevedeva un agente a due step basato su decisione LLM e
sintesi. L'implementazione corrente adotta invece un orchestratore
deterministico in \texttt{app/agent/orchestrator.py}. La scelta è più adatta a
un prototipo con modelli locali piccoli: riduce la variabilità, rende testabile
il comportamento e limita il numero di chiamate al modello.

Questa parte dell'implementazione deve quindi essere interpretata come una
soluzione esplorata e testata, non come il nucleo della soluzione finale. Le
prove condotte hanno mostrato che il modello agentico richiede più risorse
computazionali, soprattutto in termini di latenza e disponibilità di memoria,
mentre l'approccio RAG standard risulta più robusto e conveniente a parità di
hardware. Per questo motivo la componente agentica è stata mantenuta come
evidenza sperimentale e come possibile direzione futura, ma non come scelta
architetturale vincente per l'implementazione finale.

\subsection{Stato conversazionale}
\label{subsec:implementation_agent_state}

Lo stato dell'agente è rappresentato da \texttt{AgentState}. Contiene la fase
corrente, il topic della proposta, il campo del form in compilazione e una bozza
strutturata. Le fasi previste sono \texttt{triage}, \texttt{benchmark},
\texttt{ideation}, \texttt{form}, \texttt{finalize}, \texttt{search} e
\texttt{answer\_direct}. La bozza contiene i campi minimi necessari per
formalizzare una proposta: nome, descrizione, rischi, risorse e risultati
attesi. Il modello Pydantic dello stato è riportato in Appendice~\ref{app:B},
Listato~\ref{lst:agent_state}.

\begin{table}[ht]
  \centering
  \caption{Fasi del workflow agentico e comportamento associato.}
  \label{tab:agentic_workflow_phases}
  \small
  \begin{tabular}{p{3.2cm}p{4.5cm}p{5.3cm}}
    \toprule
    \textbf{Fase} & \textbf{Attivazione} & \textbf{Comportamento} \\
    \midrule
    \texttt{triage}
      & Primo messaggio o stato non definito
      & Classifica l'intento dell'utente e sceglie il ramo successivo. \\
    \texttt{benchmark}
      & Richiesta di proposta su un tema
      & Cerca Gilde precedenti utili come confronto iniziale. \\
    \texttt{ideation}
      & Benchmark completato
      & Propone una direzione progettuale e prepara la compilazione. \\
    \texttt{form}
      & Raccolta dei dati della proposta
      & Valida progressivamente nome, descrizione, rischi, risorse e output
        attesi. \\
    \texttt{finalize}
      & Campi minimi compilati
      & Restituisce una bozza coerente con il form i3Guild. \\
    \texttt{search}
      & Domanda su Gilde esistenti
      & Usa il retrieval senza avviare il flusso di proposta. \\
    \texttt{answer\_direct}
      & Messaggi semplici o saluti
      & Risponde senza interrogare Qdrant. \\
    \bottomrule
  \end{tabular}
\end{table}

Ogni risposta agentica include anche un \texttt{AgentTrace}, cioè un blocco
diagnostico con decisione, query riscritta, ragionamento sintetico, numero di
documenti recuperati e filtri estratti. Questo dato non serve al modello, ma al
frontend e alla valutazione: permette di capire perché il sistema ha scelto una
ricerca, una risposta diretta o la prosecuzione del questionario.

\subsection{Regole di orchestrazione}
\label{subsec:implementation_agent_rules}

L'orchestratore usa espressioni regolari per classificare il messaggio utente.
Messaggi di saluto vengono gestiti con risposta diretta; richieste su Gilde
esistenti attivano la ricerca; richieste di proposta o creazione avviano una
fase di benchmark, cioè una ricerca di esperienze precedenti sul topic
individuato. Quando l'utente entra nella fase di compilazione, il sistema valida
il campo corrente e passa al successivo solo se il contenuto è sufficiente.

Questa logica è meno flessibile di un agente completamente governato da LLM, ma
ha due vantaggi concreti. Primo, il comportamento è riproducibile. Secondo, il
sistema può guidare l'utente in un flusso formale senza affidare a un modello
piccolo la gestione implicita dello stato. Tuttavia, anche in questa forma più
controllata, il flusso agentico introduce complessità e overhead superiori
rispetto al percorso RAG standard. Il compromesso è quindi utile come
sperimentazione, ma non sufficiente a giustificare l'adozione nel prototipo
finale.

\subsection{Eventi di streaming agentico}
\label{subsec:implementation_agentic_stream}

La versione streaming dell'endpoint agentico restituisce eventi NDJSON. Il primo
evento comunica la decisione dell'agente; il secondo espone i dati RAG; seguono
inizio stream, token e fine stream. Anche se l'orchestratore corrente produce la
risposta completa prima di emettere i token, il contratto HTTP resta compatibile
con una futura implementazione in cui la sintesi venga generata token per token
da Ollama. Un esempio di stream agentico è riportato in
Appendice~\ref{app:B}, Listato~\ref{lst:agentic_ndjson}.

% -------------------------------------------------------------
\section{Integrazione lato Next.js}
\label{sec:implementation_nextjs}
% -------------------------------------------------------------

L'applicazione i3Guild accede al microservizio tramite
\texttt{src/lib/services/rag.service.ts}. Il client centralizza l'URL del
servizio, legge \texttt{I3GUILD\_RAG\_URL} o \texttt{RAG\_API\_URL}, serializza
le richieste e normalizza gli errori. Questa scelta evita di distribuire chiamate
HTTP dirette al servizio RAG in più componenti React o route API~\cite{nextjs2024}.

\subsection{Client di retrieval}
\label{subsec:implementation_nextjs_retrieval}

La funzione \texttt{searchRAGWithMeta()} invoca \texttt{/rag/retrieve} e
restituisce sia documenti sia metadati. In caso di errore non solleva
un'eccezione: restituisce una lista vuota e un \texttt{cutoff\_reason} pari a
\texttt{error}. Per la sola ricerca semantica questo comportamento è corretto:
il chatbot non deve bloccarsi se il servizio Python è temporaneamente non
raggiungibile. La chiamata TypeScript essenziale è riportata in
Appendice~\ref{app:B}, Listato~\ref{lst:typescript_search_rag}.

La funzione \texttt{searchRAG()} è un wrapper più semplice che restituisce solo
i documenti. Viene mantenuta per i punti dell'applicazione che non hanno bisogno
della diagnostica.

\subsection{Client chat e modalità agentica sperimentale}
\label{subsec:implementation_nextjs_chat}

Per la generazione sono disponibili \texttt{generateRAGChatReply()},
\texttt{generateRAGChatStream()} e \texttt{generateAgenticRAGChatStream()}. La
prima usa l'endpoint non streaming \texttt{/chat/reply}; le altre restituiscono
direttamente la \texttt{Response}, lasciando al chiamante la lettura incrementale
dello stream NDJSON.

La variante agentica sperimentale passa anche \texttt{agent\_state}. Questo
consente al frontend di mantenere lo stato tra turni di conversazione senza
imporre al microservizio una sessione server-side. Il servizio RAG rimane quindi
stateless: ogni richiesta contiene la cronologia e lo stato necessari per
calcolare il turno successivo. La scelta è utile per isolare la sperimentazione,
ma non modifica il flusso RAG standard adottato come riferimento finale.

% -------------------------------------------------------------
\section{Generazione della bozza di Gilda}
\label{sec:implementation_guild_draft}
% -------------------------------------------------------------

Accanto alla chat RAG, i3Guild include un flusso per generare una bozza di
Gilda a partire dalla conversazione. La funzionalità è documentata come
\emph{Guild Draft Generation} e si attiva dal pulsante "Proponi bozza" nella
chat. Il flusso non crea direttamente una nuova Gilda e non invia l'utente alla
pagina di creazione: produce invece un oggetto strutturato che rappresenta un
draft, modificabile e verificabile prima di qualunque salvataggio applicativo.

Il formato atteso è \texttt{GuildDraft}. Include nome, sommario, obiettivi,
risultati attesi, rischi, opportunità, budget, risorse, journal, modalità di
partecipazione, numero minimo e massimo di partecipanti, indicazione di Gilda
individuale e avatar. Il servizio costruisce un prompt istruttivo a partire
dalla conversazione, invoca Ollama e prova a estrarre dalla risposta un JSON
compatibile con questo schema. La validazione è quindi parte del contratto
applicativo: una risposta testuale plausibile, ma non riconducibile allo schema
atteso, non viene trattata come bozza valida.

La gestione degli errori è intenzionalmente esplicita. Se la conversazione non è
sufficiente, se il modello restituisce JSON non valido o se mancano campi
obbligatori, il sistema informa l'utente del problema e viene invitato a riprovare, 
eventualmente aggiungendo dettagli più chiari nella conversazione. Questo 
comportamento rende misurabile il tasso di fallimento della generazione. Nella
valutazione sperimentale, tale tasso può essere espresso come rapporto tra
bozze non valide e numero totale di richieste di generazione, duratne i test 
effettuati il tasso di errori è stato di: $n_{fallimenti} / N_{test}$. \textbf{TEST DA FARE}


% -------------------------------------------------------------
\section{Aspetti operativi}
\label{sec:implementation_operations}
% -------------------------------------------------------------

L'implementazione include alcune scelte operative che non cambiano la logica RAG
ma ne migliorano l'affidabilità durante l'uso locale.

\subsection{Timeout e dipendenze lente}
\label{subsec:implementation_timeouts}

Ollama può impiegare diversi secondi, o più di un minuto su CPU, per caricare un
modello non ancora residente in memoria. Il problema si presenta soprattutto
alla prima richiesta dopo l'avvio del container: il modello è installato, ma il
processo deve ancora caricarlo e inizializzare l'inferenza. Per questo il
servizio Python usa timeout più lunghi nelle chiamate streaming verso Ollama.

La stessa attenzione serve lato Next.js. Le chiamate server-side effettuate con
\texttt{fetch} nel runtime Node.js usano Undici come client HTTP sottostante.
Configurare un \texttt{Agent} di Undici significa impostare il dispatcher usato
da quelle richieste, estendendo parametri come \texttt{headersTimeout} e
\texttt{bodyTimeout}. In questo modo il frontend non interrompe prematuramente
la richiesta mentre il microservizio RAG sta ancora attendendo la prima risposta
dal modello locale.

\subsection{System prompt e modello derivato}
\label{subsec:implementation_prompt_cache}

Una seconda ottimizzazione operativa riguarda la gestione del system prompt.
Nella configurazione standard il prompt di sistema viene inviato a ogni
richiesta insieme alla conversazione e al contesto recuperato. Per ridurre la
duplicazione del payload, il progetto considera anche la creazione di un modello
Ollama derivato tramite l'endpoint \texttt{/api/create}, incorporando il system
prompt nella definizione del modello. Il nome del modello derivato include un
hash del prompt, così una modifica del prompt produce una nuova configurazione
senza sovrascrivere quella precedente.

Questa soluzione migliora soprattutto la stabilità configurativa e riduce la
quantità di testo serializzata nelle richieste. Non va però interpretata come un
riuso effettivo della KV Cache del modello: nei benchmark registrati il payload
JSON si riduce in modo marcato, ma il numero di token valutati da Ollama rimane
sostanzialmente invariato. L'ottimizzazione è quindi utile per rendere più
pulita e riproducibile la configurazione, non per eliminare il costo computazionale
del prompt durante l'inferenza. \textbf{TEST DA FARE}

\subsection{Tracciabilità}
\label{subsec:implementation_tracing}

Il servizio espone dati diagnostici in più punti: \texttt{cutoff\_reason},
\texttt{topic\_query}, \texttt{extracted\_filters}, warning RAM e
\texttt{AgentTrace}. Questi campi non sono solo utili per il debug. Permettono
di ricostruire il comportamento del sistema durante la valutazione sperimentale:
quale query è stata effettivamente cercata, quanti documenti sono stati
recuperati, perché alcuni risultati sono stati scartati e quale fase agentica ha
prodotto la risposta.

% -------------------------------------------------------------
\section{Sintesi del capitolo}
\label{sec:implementation_summary}
% -------------------------------------------------------------

Il sistema implementato combina componenti semplici, ma separati con precisione.
La pipeline di embedding prepara l'indice a partire da PostgreSQL e da sorgenti
opzionali; il microservizio RAG espone retrieval e generazione; Qdrant gestisce
la ricerca vettoriale; Ollama mantiene l'inferenza in locale; Next.js integra le
funzioni nel flusso utente. Le parti più specifiche del lavoro sono la
normalizzazione del corpus i3Guild, il retrieval dinamico con query rewriting e
deduplicazione. Il workflow agentico deterministico è stato invece trattato come
una linea sperimentale: ha permesso di valutare un flusso più guidato per la
proposta di nuove Gilde, ma non è risultato preferibile all'approccio classico
nelle condizioni hardware disponibili. Queste scelte rendono il prototipo
coerente con i vincoli di privacy e di risorse discussi nei capitoli precedenti,
lasciando aperti margini di evoluzione su sincronizzazione incrementale, agenti
più flessibili e valutazione sistematica della qualità delle risposte.
